\section*{摘要}
本文围绕 Lorenz96 数据同化中的一个具体问题展开：在经典 LETKF 已经较为稳定的前提下，量子信息是否还能以可解释、可复现且资源受限的方式提供额外增益。为此，本文不采用全量子替代路线，而是保留经典 LETKF 作为主骨架，将量子模块限制为局地结构信息提供者。具体做法是：利用 $4$ qubits 量子特征映射构造局地状态窗口与观测窗口之间的量子相似度，用其对经典距离先验进行弱融合，并在集合扰动方向上引入由局地耦合系数控制的弱 $J$ 修正。本文在固定实验设置下，对 \texttt{fixed}、\texttt{corr} 和四组 $(\lambda_\rho,\lambda_J)$ 量子弱融合配置进行了受控比较。结果显示，最优配置 $(\lambda_\rho,\lambda_J)=(0.02,0.003)$ 在 test\_1 上取得 $0.12698$ 的 RMSE，相比 \texttt{fixed} 的 $0.12873$ 下降约 $1.35\%$，同时优于 \texttt{corr} 的 $0.14279$。进一步的误差分布、逐时间步趋势、逐维度改善和代表性状态切片分析表明，该增益主要体现为局地误差尾部收缩和部分难点维度的恢复改善。基于当前单数据设定的证据，本文支持如下较谨慎的结论：量子相似度更适合作为局地加权与弱方向修正信号，而不是直接主导整体分析更新；这一判断适用于本文的受控实验条件，尚需在更大规模数据和更多方法变体上进一步验证。
